\documentclass[12pt,aspectratio=169]{beamer}

\usetheme{Madrid}
\usecolortheme{whale}
\useinnertheme{rectangles}
\setbeamertemplate{navigation symbols}{}
\usepackage{fontspec}
\usepackage{xeCJK}
\setCJKmainfont{標楷體}
\setCJKsansfont{標楷體}
\setCJKmonofont{標楷體}

\usepackage{booktabs}
\usepackage{tikz}
\usetikzlibrary{shapes.geometric, arrows.meta, positioning}
\usepackage{listings}
\usepackage{xcolor}

% Custom palette: deep navy header, clean body contrast
\definecolor{navydark}{RGB}{13,38,75}
\definecolor{navymid}{RGB}{25,65,120}
\definecolor{navylight}{RGB}{200,215,240}
\setbeamercolor{palette primary}{bg=navydark,fg=white}
\setbeamercolor{palette secondary}{bg=navymid,fg=white}
\setbeamercolor{palette tertiary}{bg=navylight,fg=navydark}
\setbeamercolor{frametitle}{bg=navydark,fg=white}
\setbeamercolor{title}{fg=white}
\setbeamercolor{subtitle}{fg=navylight}
\setbeamercolor{normal text}{fg=black!90,bg=white}
\setbeamercolor{item}{fg=navymid}
\setbeamercolor{subitem}{fg=navymid!70}

% Block colors
\setbeamercolor{block title}{bg=navymid,fg=white}
\setbeamercolor{block body}{bg=navylight!40,fg=black!90}
\setbeamercolor{block title alerted}{bg=red!65!black,fg=white}
\setbeamercolor{block body alerted}{bg=red!8,fg=black!90}
\setbeamercolor{block title example}{bg=teal!60!black,fg=white}
\setbeamercolor{block body example}{bg=teal!8,fg=black!90}

% Code block colors (Dracula-inspired, high contrast)
\definecolor{codebg}{RGB}{30,32,43}
\definecolor{codefg}{RGB}{248,248,242}
\definecolor{codekey}{RGB}{255,121,198}
\definecolor{codestr}{RGB}{241,250,140}
\definecolor{codecomment}{RGB}{98,114,164}

\lstdefinestyle{dark}{
  backgroundcolor=\color{codebg},
  basicstyle=\ttfamily\small\color{codefg},
  keywordstyle=\color{codekey}\bfseries,
  stringstyle=\color{codestr},
  commentstyle=\color{codecomment}\itshape,
  breaklines=true,
  showstringspaces=false,
  frame=single,
  rulecolor=\color{navymid},
  framerule=1pt,
  columns=fullflexible,
  xleftmargin=4pt,
  xrightmargin=4pt,
}
\setbeamerfont{frametitle}{size=\Large,series=\bfseries}
\setbeamerfont{title}{size=\LARGE,series=\bfseries}
\setbeamerfont{itemize/enumerate body}{size=\normalsize}
\setbeamerfont{itemize/enumerate subbody}{size=\small}
\lstset{style=dark}

% Mini compromise chain for sidebar use.
% Args 1-6: style per node (use "mcbox" normal or "mchl" highlighted).
% Nodes in order: Caldera, splunkd, powershell, Akagi64, slui, cmd
\newcommand{\minichain}[6]{%
\begin{tikzpicture}[
  node distance=0.12cm,
  mcbox/.style={rectangle, rounded corners, fill=blue!50!black, text=white,
                minimum width=3.2cm, minimum height=0.5cm, font=\scriptsize, inner sep=2pt},
  mchl/.style={rectangle, rounded corners, fill=red!70!black, text=white,
               draw=yellow, line width=1pt,
               minimum width=3.2cm, minimum height=0.5cm, font=\scriptsize\bfseries, inner sep=2pt},
  arr/.style={-{Stealth[length=1.2mm]}, thin, draw=gray}
]
\node[font=\scriptsize\bfseries] (title) {Compromise Chain};
\node[#1, below=0.2cm of title] (s) {Caldera Server};
\node[#2, below=of s]  (a)  {splunkd.exe (M)};
\node[#3, below=of a]  (p)  {powershell.exe (M)};
\node[#4, below=of p]  (k)  {Akagi64.exe (M)};
\node[#5, below=of k]  (sl) {slui.exe (H)};
\node[#6, below=of sl] (c)  {cmd.exe (H)};
\draw[arr] (s) -- (a); \draw[arr] (a) -- (p); \draw[arr] (p) -- (k);
\draw[arr] (k) -- (sl); \draw[arr] (sl) -- (c);
\end{tikzpicture}}

\title{MITRE Caldera 攻擊模擬與日誌分析}
\subtitle{T1548.002 - Slui File Handler Hijack}
\author{趙子佾, 吳秉哲}

\begin{document}

\maketitle
% ------------------------------------------------------------
\begin{frame}{攻擊技術說明 T1548.002}
\begin{itemize}
  \item \textbf{Tactic}：Privilege Escalation（權限提升）／ Defense Evasion
  \item \textbf{技術名稱}：Abuse Elevation Control Mechanism - Bypass User Account Control
  \item \textbf{子技術}：Slui File Handler Hijack
  \item \textbf{原理}：
  \begin{itemize}
    \item 先改寫 \texttt{HKCU\textbackslash Software\textbackslash Classes\textbackslash exefile} File Handler，再呼叫 \texttt{slui.exe}
    \item \texttt{slui.exe} 啟動時讀到被竄改的 HKCU，且本身具有 \texttt{autoElevate=true}
    \item 以 High Integrity 啟動被劫持的 \texttt{cmd.exe}
    \item 從 Medium Integrity $\rightarrow$ High Integrity，\textbf{不需使用者互動，不彈 UAC 視窗}
  \end{itemize}
  \item \textbf{工具}：UACME（\texttt{Akagi64.exe}）Method 45 — \textbf{開源} GitHub 專案（hfiref0x/UACME），收錄 90+ 種 UAC bypass 手法
\end{itemize}
\end{frame}

% ------------------------------------------------------------
\begin{frame}{背景知識 (1)：Manifest 與 autoElevate}
\begin{itemize}
  \item \textbf{Windows Manifest}：附加在 \texttt{.exe} 內的 XML 設定（\texttt{<assembly>}），\textbf{所有 Windows 執行檔都可使用}，並非 \texttt{slui.exe} 專屬。常見用途：宣告相依 DLL、UI 主題、DPI 感知、UAC 行為。
  \item \textbf{\texttt{autoElevate=true}}：Manifest 內 \texttt{<requestedExecutionLevel>} 旁的旗標，告訴系統「此程式為 Microsoft 信任程式，可在符合條件時\textbf{自動升至 High Integrity}」。
  \item \textbf{自動放行條件（同時滿足）}：
  \begin{enumerate}
    \item Manifest 含 \texttt{autoElevate=true}
    \item 位於 \texttt{C:\textbackslash Windows\textbackslash System32\textbackslash}（受信任目錄）
    \item 具有 Microsoft 數位簽章
  \end{enumerate}
  \item \textbf{一般使用者可否呼叫 slui.exe？} 可以，它本身就是普通可執行檔，任何使用者皆可啟動，啟動時系統依上述條件自動賦予 High Integrity。
\end{itemize}
\end{frame}

% ------------------------------------------------------------
\begin{frame}[fragile]{背景知識 (2)：UAC autoElevate 白名單}
Windows 預設的 autoElevate 程式（節錄，皆位於 \texttt{System32}）：
\begin{lstlisting}[style=dark,basicstyle=\scriptsize\ttfamily\color{codefg}]
slui.exe              ComputerDefaults.exe   fodhelper.exe
sdclt.exe             WSReset.exe            CompMgmtLauncher.exe
changepk.exe          eventvwr.exe (legacy)  DeviceEject.exe
EASPolicyManagerBrokerHost.exe                msconfig.exe
\end{lstlisting}
\textbf{取得方式}：以 PowerShell 掃描 \texttt{System32} 內所有 \texttt{.exe}，抽取 manifest 資源並比對 \texttt{autoElevate} 標籤。
\begin{lstlisting}[style=dark,basicstyle=\scriptsize\ttfamily\color{codefg}]
Get-ChildItem C:\Windows\System32\*.exe | ForEach-Object {
  $m = (Select-String -Path $_.FullName -Pattern "autoElevate" `
        -Encoding unicode -SimpleMatch -ErrorAction SilentlyContinue)
  if ($m) { $_.Name }
}
\end{lstlisting}
這些每一支都是潛在的 UAC bypass 切入點 — UACME 之所以能收錄 90+ 種手法，正是因為白名單上有許多程式都會讀 HKCU。
\end{frame}

% ------------------------------------------------------------
\begin{frame}{背景知識 (3)：Registry、File Handler、HKCU}
\begin{itemize}
  \item \textbf{Registry（登錄檔）}：Windows 的階層式設定資料庫，所有系統／應用程式設定皆儲存於此，以「Hive $\rightarrow$ Key $\rightarrow$ Value」結構組織。
  \item \textbf{File Handler}：登錄檔中描述「某副檔名／檔案類型由哪支程式開啟」的設定。例如 \texttt{exefile} 就代表「\texttt{.exe} 檔的處理者」，其 \texttt{shell\textbackslash open\textbackslash command} 即為「執行 .exe 時要呼叫哪一條命令列」。
  \item \textbf{HKCU}（HKEY\_CURRENT\_USER）：當前登入使用者的個人設定 hive，\textbf{使用者本人即可讀寫}，不需管理員權限。
  \item \textbf{HKLM}（HKEY\_LOCAL\_MACHINE）：全機器設定，僅管理員可寫。
  \item \textbf{查詢順序（COM/Shell API 規定）}：HKCU \textbf{優先} $\rightarrow$ 找不到才 fallback 到 HKLM。\\
  $\Rightarrow$ 若 HKCU 寫了一份「假的」exefile handler，整支 slui.exe 看到的就是假的版本。
\end{itemize}
\end{frame}

% ------------------------------------------------------------
\begin{frame}{背景知識 (4)：HKCU 為何能被普通使用者寫入？}
\textbf{設計理念}：HKCU 的設計初衷是「讓使用者可在不影響其他人的前提下，自訂自己的應用程式行為」（如預設瀏覽器、檔案關聯、佈景主題）。因此 HKCU 對使用者本人開放完整讀寫權限。
\begin{itemize}
  \item HKCU 優先於 HKLM 是\textbf{功能}（讓使用者自訂）而非錯誤
  \item 但這對「會以高權限執行的程式」而言，等於信任了一份低權限可寫的設定
\end{itemize}

\textbf{為何至今未修掉？}
\begin{itemize}
  \item 修掉「HKCU 優先」會破壞海量現有應用程式的檔案關聯邏輯（向後相容性災難）
  \item 微軟採取的折衷做法：在新版 Windows 中，將「會 autoElevate 的關鍵程式」逐支從讀 HKCU 改為只讀 HKLM，或改用其他驗證機制（例如 \texttt{fodhelper}、\texttt{eventvwr} 等手法已在較新版本被修補）
  \item \textbf{slui.exe 的此手法在某些 Windows 11 版本仍可使用}，因為微軟採「逐支修補」而非「全面禁止」
\end{itemize}
\end{frame}

% ------------------------------------------------------------
\begin{frame}[fragile]{slui.exe 是什麼？為何擁有最高權限？}
\textbf{slui.exe — Windows Software Licensing UI}
\begin{itemize}
  \item 位置：\texttt{C:\textbackslash Windows\textbackslash System32\textbackslash slui.exe}
  \item Microsoft 簽章，負責 Windows 啟用／授權 UI
\end{itemize}

\vspace{0.4em}
\textbf{為何自動以 High Integrity 執行？}
\begin{itemize}
  \item Manifest 含 \texttt{autoElevate=true}
  \item 位於 System32、有 Microsoft 簽章 $\rightarrow$ 在 UAC 白名單(不會跳UAC確認)
  \item consent.exe 會負責確認程式是否被授權。若沒有，則會 負責在執行需管理員權限的操作時彈出授權視窗
  \item 符合條件者 consent.exe \textbf{自動放行}，不顯示 UAC 彈窗
\end{itemize}

\vspace{0.3em}
\textbf{關鍵弱點}：slui.exe 會查詢 HKCU 下的 \texttt{exefile} File Handler，而 HKCU 普通使用者即可寫入 $\rightarrow$ 可被 hijack。

\vspace{0.3em}
\textbf{一般使用者可否呼叫 slui.exe？} 可以——它是普通可執行檔，無存取限制；正因如此攻擊者才能在 Medium Integrity 下主動觸發它。
\end{frame}

% ------------------------------------------------------------
\begin{frame}{攻擊執行流程（詳細版）}
\centering
\begin{tikzpicture}[
  node distance=0.28cm and 0.5cm,
  box/.style={rectangle, rounded corners, draw=white, fill=blue!60!black,
              text=white, minimum width=6.2cm, minimum height=0.55cm, font=\scriptsize},
  hbox/.style={rectangle, rounded corners, draw=white, fill=red!70!black,
               text=white, minimum width=6.2cm, minimum height=0.55cm, font=\scriptsize},
  reg/.style={rectangle, rounded corners, draw=white, fill=orange!80!black,
              text=white, minimum width=5.4cm, minimum height=0.5cm, font=\scriptsize},
  branch/.style={rectangle, rounded corners, draw=orange!80!black, line width=0.8pt,
                 fill=orange!10, text=black, minimum width=5.6cm, minimum height=0.5cm,
                 font=\scriptsize, align=left},
  arr/.style={-{Stealth[length=2.2mm]}, thick, draw=white!60!black},
  barr/.style={-{Stealth[length=1.8mm]}, thin, draw=orange!70!black, dashed}
]
% --- Main chain (Sysmon view) ---
\node[box] (s)  {1. Caldera Server 下發 ability};
\node[box, below=of s]  (a)  {2. splunkd.exe (Medium) [01:09:44]};
\node[box, below=of a]  (p)  {3. powershell.exe -ExecutionPolicy Bypass (Medium)};
\node[box, below=of p]  (k)  {4. Akagi64.exe 45 cmd.exe (Medium) [01:09:44]};
\node[reg, below=of k]  (r)  {5. Registry Hijack 完成（見右側 Procmon 時間線）};
\node[box, below=of r]  (q)  {6. Akagi64 觸發 slui.exe};
\node[hbox, below=of q] (sl) {7. slui.exe autoElevate $\rightarrow$ \textbf{High}，讀到假 HKCU [01:09:45]};
\node[hbox, below=of sl] (c) {8. cmd.exe (\textbf{High}) [01:09:46] $\leftarrow$ 攻擊成功};

\draw[arr] (s) -- (a); \draw[arr] (a) -- (p); \draw[arr] (p) -- (k);
\draw[arr] (k) -- (r); \draw[arr] (r) -- (q);
\draw[arr] (q) -- (sl); \draw[arr] (sl) -- (c);

% --- Procmon branch: registry writes timeline (right of step 4-5) ---
\node[font=\scriptsize\bfseries, right=1.3cm of k, anchor=west,
      text=orange!80!black] (btitle) {Procmon Registry 時間線};
\node[branch, below=0.15cm of btitle.south west, anchor=north west] (b1)
   {01:09:44.952 \textbf{RegCreateKey}\\HKCU\textbackslash...\textbackslash exefile\textbackslash shell\textbackslash open\textbackslash command};
\node[branch, below=of b1] (b2)
   {01:09:44.953 \textbf{RegSetValue} (REG\_SZ)\\Data = C:\textbackslash Windows\textbackslash System32\textbackslash cmd.exe};
\node[branch, below=of b2] (b3)
   {01:09:44.958 \textbf{QueryDirectory} slui.exe\\（確認觸發目標存在）};

\draw[barr] (k.east) to[out=0,in=180] (b1.west);
\draw[barr] (b1) -- (b2);
\draw[barr] (b2) -- (b3);
\draw[barr] (b3.west) to[out=180,in=0] (r.east);

% legend
\node[font=\tiny, below=0.15cm of c, anchor=north]
   {主鏈：Sysmon 視角（行程建立）｜分支：Procmon 視角（Registry 寫入細節）};
\end{tikzpicture}
\end{frame}

% ------------------------------------------------------------
\begin{frame}[fragile]{slui.exe 與登錄檔的關係}
\textbf{slui.exe 為何會去讀登錄檔？}
\begin{itemize}
  \item 內部會\textbf{開啟 \texttt{.exe}} $\rightarrow$ 經由 \texttt{exefile} ProgID 查「誰來開啟」
  \item 查詢順序：\textcolor{red!70!black}{\textbf{HKCU 優先}} $\rightarrow$ HKLM（HKCU 普通使用者可寫）
\end{itemize}

\vspace{0.3em}
\textbf{關鍵登錄檔 Key：}\\
\texttt{\small HKCU\textbackslash Software\textbackslash Classes\textbackslash exefile}

\vspace{0.3em}
\textbf{攻擊正確順序（四步）}：
\begin{enumerate}
  \item Akagi64.exe（Medium）\textbf{先}寫入 HKCU exefile handler $\rightarrow$ 指向 \texttt{cmd.exe}
  \item Akagi64.exe \textbf{後}才啟動 \texttt{slui.exe}（autoElevate $\rightarrow$ 系統自動升為 High）
  \item slui.exe（High）讀 HKCU exefile handler $\rightarrow$ 命中攻擊者預先植入的 \texttt{cmd.exe}
  \item slui.exe 以 High 建立子行程 $\rightarrow$ \texttt{cmd.exe} \textbf{繼承 High}
\end{enumerate}

\begin{alertblock}{核心原理}
\small
cmd.exe 由「已是 High」的 slui.exe 啟動，因此一出生就是 High。\textbf{權限不是被提升，而是被「借用」}。
\end{alertblock}
\end{frame}

% ------------------------------------------------------------
\begin{frame}[fragile]{原始紀錄檔規模與篩選方式}
\begin{columns}[T]
\begin{column}{0.5\textwidth}
\textbf{\small 原始日誌規模}
\begin{table}
\centering
\footnotesize
\setlength{\tabcolsep}{4pt}
\begin{tabular}{@{}llrr@{}}
\toprule
檔案 & 來源 & 大小 & 筆數 \\
\midrule
\texttt{sysmon.log}  & Sysmon EVTX   & 7.3 KB & 150 \\
\texttt{procmon.CSV} & Procmon Trace & 3.9 MB & 23{,}938 \\
\bottomrule
\end{tabular}
\end{table}
\end{column}
\begin{column}{0.5\textwidth}
\textbf{\small Procmon 篩選（GUI Filter）}
\begin{lstlisting}[style=dark,basicstyle=\scriptsize\ttfamily\color{codefg}]
Column   : Process Name
Relation : is
Value    : Akagi64.exe
Action   : Include
\end{lstlisting}
\end{column}
\end{columns}

\vspace{-0.5em}
\textbf{\small Sysmon 篩選（PowerShell 讀 EVTX，鎖定時間區間與攻擊鏈關鍵字）}
\begin{lstlisting}[style=dark,basicstyle=\scriptsize\ttfamily\color{codefg}]
Get-WinEvent -LogName "Microsoft-Windows-Sysmon/Operational" |
  Where-Object {
    $_.TimeCreated -gt [datetime]"2026-04-22 01:10:44" -and
    $_.TimeCreated -lt [datetime]"2026-04-22 01:12:44" -and
    ($_.Message -like "*Akagi*" -or
     $_.Message -like "*slui*"  -or
     $_.Message -like "*sandcat*")
  } | Select-Object TimeCreated, Id, Message | Format-List |
  Out-File "C:\Users\Public\sysmon_log.txt" -Encoding UTF8
\end{lstlisting}
\end{frame}

% ------------------------------------------------------------
\begin{frame}[fragile]{Sysmon 日誌分析 - Stage 1: Agent 啟動 (01:09:44)}
\begin{columns}[T]
\begin{column}{0.65\textwidth}
\begin{table}
\centering
\scriptsize
\begin{tabular}{@{}lll@{}}
\toprule
時間 & Process (PID) & IntegrityLevel \\
\midrule
01:09:44 & splunkd.exe (4008)  & Medium \\
01:09:44 & powershell.exe (5552)  & Medium \\
\bottomrule
\end{tabular}
\end{table}

\vspace{0.3em}
\textbf{\footnotesize 攻擊啟動：}
\begin{lstlisting}[style=dark,basicstyle=\scriptsize\ttfamily\color{codefg}]
Image: C:\Users\Public\splunkd.exe
Server: http://10.1.106.114:8888
-group red
\end{lstlisting}

\vspace{0.2em}
\textbf{\footnotesize PowerShell 執行：}
\begin{lstlisting}[style=dark,basicstyle=\scriptsize\ttfamily\color{codefg}]
Image: powershell.exe (5552)
CommandLine: -ExecutionPolicy Bypass
  -C "./Akagi64.exe 45 ..."
\end{lstlisting}
\end{column}
\begin{column}{0.35\textwidth}
\centering
\minichain{mchl}{mchl}{mcbox}{mcbox}{mcbox}{mcbox}
\end{column}
\end{columns}
\end{frame}

% ------------------------------------------------------------
\begin{frame}[fragile]{Sysmon 日誌分析 - Stage 2: UAC Bypass (01:09:44)}
\begin{columns}[T]
\begin{column}{0.65\textwidth}
\begin{table}
\centering
\scriptsize
\begin{tabular}{@{}lll@{}}
\toprule
時間 & Process (PID) & IntegrityLevel \\
\midrule
01:09:44 & powershell.exe (5552) & Medium \\
01:09:44 & Akagi64.exe (11604) & Medium \\
\bottomrule
\end{tabular}
\end{table}

\vspace{0.3em}
\textbf{\footnotesize UAC Bypass 工具啟動 (Method 45)：}
\begin{lstlisting}[style=dark,basicstyle=\scriptsize\ttfamily\color{codefg}]
Image: Akagi64.exe
ProcessId: 11604
CommandLine: Akagi64.exe 45
  C:\Windows\System32\cmd.exe
Description: UACMe main module
FileVersion: 3.2.2.1911
IntegrityLevel: Medium
\end{lstlisting}
\end{column}
\begin{column}{0.35\textwidth}
\centering
\minichain{mcbox}{mchl}{mchl}{mcbox}{mcbox}{mcbox}
\end{column}
\end{columns}
\end{frame}

% ------------------------------------------------------------
\begin{frame}[fragile]{Sysmon 日誌分析 - Stage 3: Registry Hijack (01:09:45)}
\begin{columns}[T]
\begin{column}{0.65\textwidth}
\begin{table}
\centering
\scriptsize
\begin{tabular}{@{}lll@{}}
\toprule
時間 & Process (PID) & IntegrityLevel \\
\midrule
01:09:44 & Akagi64.exe (11604) & Medium \\
01:09:45 & slui.exe (3272) & \textbf{High} \\
01:09:45 & slui.exe 0x03 (3436) & \textbf{High} \\
\bottomrule
\end{tabular}
\end{table}

\vspace{0.3em}
\textbf{\footnotesize slui.exe 連鎖啟動（Registry Hijack）：}
\begin{lstlisting}[style=dark,basicstyle=\scriptsize\ttfamily\color{codefg}]
Image: slui.exe (PID 3272)
CommandLine: slui.exe
ParentImage: Akagi64.exe (11604)
IntegrityLevel: High ← Auto-elevated!

→ slui.exe (PID 3436) 0x03
CommandLine: slui.exe 0x03
IntegrityLevel: High
\end{lstlisting}
\end{column}
\begin{column}{0.35\textwidth}
\centering
\minichain{mcbox}{mcbox}{mcbox}{mchl}{mchl}{mcbox}
\end{column}
\end{columns}
\end{frame}

% ------------------------------------------------------------
\begin{frame}[fragile]{Sysmon 日誌分析 - Stage 4: 權限繼承 (01:09:46)}
\begin{columns}[T]
\begin{column}{0.65\textwidth}
\begin{table}
\centering
\scriptsize
\begin{tabular}{@{}lll@{}}
\toprule
時間 & Process (PID) & IntegrityLevel \\
\midrule
01:09:45 & slui.exe 0x03 (3436) & \textbf{High} \\
01:09:46 & cmd.exe (4084) & \textbf{High} \\
01:09:50 & Akagi64.exe & 終止 \\
\bottomrule
\end{tabular}
\end{table}

\vspace{0.3em}
\textbf{\footnotesize 最終結果：High Integrity Shell}
\begin{lstlisting}[style=dark,basicstyle=\scriptsize\ttfamily\color{codefg}]
Image: cmd.exe
ProcessId: 4084
ParentProcessId: 3436
ParentImage: slui.exe
IntegrityLevel: High [SUCCESS!]
Description: Windows Command Processor
FileVersion: 10.0.26100.8115
MD5: 77F0062F490BCC702...
\end{lstlisting}
\end{column}
\begin{column}{0.35\textwidth}
\centering
\minichain{mcbox}{mcbox}{mcbox}{mcbox}{mchl}{mchl}
\end{column}
\end{columns}
\end{frame}


% ------------------------------------------------------------
% \begin{frame}[fragile]{Sysmon 日誌分析 - 關鍵偵測點}
% \begin{alertblock}{\large 攻擊特徵（IOC）}
% \footnotesize
% \begin{itemize}
%   \item \textbf{異常進程路徑}：\texttt{splunkd.exe}、\texttt{Akagi64.exe} 位於 \texttt{C:\textbackslash Users\textbackslash Public\textbackslash}
%   \item \textbf{可疑父子關係}：
%   \begin{itemize}
%     \item \texttt{splunkd.exe} → \texttt{powershell.exe} → \texttt{Akagi64.exe}（連續提升權限工具）
%     \item \texttt{slui.exe} 作為 \texttt{cmd.exe} 的父行程（異常）
%   \end{itemize}
%   \item \textbf{完整性級別提升}：Medium → \textbf{High}（UAC Bypass 成功跡象）
%   \item \textbf{時間關聯性}：1秒內完成整個 UAC Bypass 鏈（自動化攻擊）
%   \item \textbf{工具指紋}：UACMe (Akagi64.exe, FileVersion 3.2.2.1911)
%   \item \textbf{Hash 值}（偵測指標）：
%   \begin{itemize}
%     \item Akagi64.exe MD5: \texttt{CA3FE69BFDBD536856BF36BBB8C949D4}
%     \item cmd.exe MD5: \texttt{77F0062F490BCC7023763A422E561945}
%   \end{itemize}
% \end{itemize}
% \end{alertblock}
% \end{frame}

% ------------------------------------------------------------
\begin{frame}[fragile]{Procmon 輔助分析 - Registry Hijack}
Procmon 補充 Sysmon 所看不到的 Registry 操作，揭露 Method 45 核心手法：
\begin{lstlisting}[style=dark,basicstyle=\scriptsize\ttfamily\color{codefg}]
01:09:44.952  Akagi64.exe  RegCreateKey
   HKCU/Software/Classes/exefile/shell/open/command

01:09:44.953  Akagi64.exe  RegSetValue
   HKCU/Software/Classes/exefile/shell/open/command
   Type: REG_SZ
   Data: C:\Windows\System32\cmd.exe

01:09:44.958  Akagi64.exe  QueryDirectory
   slui.exe
\end{lstlisting}

\textbf{手法說明}：攻擊者將 \texttt{exefile} 的 shell open command 改寫為 \texttt{cmd.exe}，
接著觸發具有 \texttt{autoElevate} 屬性的 \texttt{slui.exe}，
由其以 High Integrity 啟動被劫持的 \texttt{cmd.exe}。

\vspace{0.3em}
\textbf{細節補充：}
\begin{itemize}
  \footnotesize
  \item \texttt{REG\_SZ}：Registry 中「以 null 結尾的 Unicode 字串」型別。是\textbf{覆寫}（若原本沒有 key 則 \texttt{RegCreateKey} 建立後 \texttt{RegSetValue} 寫入；若原本有，則整個值被取代），不是 null 寫入。
  \item \texttt{QueryDirectory}：Akagi64 在列舉資料夾項目（如確認 \texttt{slui.exe} 路徑、檢查 System32 內檔案是否存在），屬於攻擊前的探測動作。
\end{itemize}

\vspace{0.2em}
\textbf{為何沒出現「輸入金鑰」UI 而是直接跳 cmd？}\\
\footnotesize
\texttt{slui.exe 0x03} 內部會\textbf{以 ShellExecute 開啟一支 .exe} 來顯示金鑰輸入介面。系統查 exefile handler 時 HKCU 已被改成 \texttt{cmd.exe}，因此 slui 不是執行原本的金鑰 UI，而是\textbf{以 High Integrity 啟動 cmd.exe}，原本的金鑰輸入流程因此被偷天換日。
\end{frame}

% ------------------------------------------------------------
% \begin{frame}{結論}
% \begin{itemize}
%   \item 成功使用 MITRE Caldera 模擬 \textbf{T1548.002} Slui File Handler Hijack 攻擊
%   \item Sysmon 完整記錄從 Agent 啟動到 UAC Bypass 的完整攻擊鏈
%   \item 攻擊者使用 \texttt{splunkd.exe} 偽裝成合法程式，並透過 \texttt{Akagi64.exe} 執行 Method 45 繞過 UAC
%   \item Procmon 補足 Registry 層級證據：\texttt{HKCU\textbackslash Software\textbackslash Classes\textbackslash exefile\textbackslash shell\textbackslash open\textbackslash command} 被劫持
%   \item \textbf{IntegrityLevel 從 Medium 提升至 High} 是關鍵偵測指標
% \end{itemize}
% \end{frame}

% ------------------------------------------------------------
\begin{frame}
\centering
\vfill
{\Huge\bfseries Thank You}\\[1em]
{\Large Q \& A}
\vfill
\end{frame}

\end{document}
